% ========================================================
%  《习近平新时代中国特色社会主义思想概论》  全真模拟试卷（含参考答案）
%  覆盖范围：导论至第十七章，重点高频考点 + 客观题细节
% ========================================================
\documentclass[12pt, a4paper]{article}

% ---- 中文支持 ----
\usepackage[UTF8]{ctex}

% ---- 数学 ----
\usepackage{amsmath, amssymb, amsthm}
\usepackage{mathtools}

% ---- 页面布局 ----
\usepackage[top=2.5cm, bottom=2.5cm, left=2.8cm, right=2.8cm]{geometry}

% ---- 表格 ----
\usepackage{booktabs, array, multirow, makecell}
\usepackage{tabularx}

% ---- 页眉页脚 ----
\usepackage{fancyhdr}
\usepackage{lastpage}
\pagestyle{fancy}
\fancyhf{}
\renewcommand{\headrulewidth}{0.6pt}
\lhead{\small 习近平新时代中国特色社会主义思想概论·全真模拟}
\rhead{\small 共 \pageref{LastPage} 页}
\cfoot{\small 第 \thepage\ 页}

% ---- 计数器与样式 ----
\usepackage{enumitem}
\usepackage{xcolor}
\usepackage{tcolorbox}
\tcbuselibrary{skins, breakable}

% ---- 填空线 ----
\newcommand{\blank}[1]{\underline{\hspace{#1}}}

% ---- 答案盒子 ----
\newtcolorbox{answerbox}[1][]{
  breakable,
  colback=gray!6,
  colframe=gray!50,
  fonttitle=\bfseries\small,
  title=参考答案,
  #1
}

% ---- 题号格式 ----
\newcounter{bigq}
\newcommand{\BigQ}[1]{%
  \stepcounter{bigq}%
  \vspace{6pt}%
  \noindent{\large\bfseries 第\chinese{bigq}大题\quad #1}%
  \vspace{4pt}\par
}

\newcounter{subq}[bigq]
\newcommand{\SubQ}{%
  \stepcounter{subq}%
  \noindent\textbf{（\arabic{subq}）}\quad
}

\newcounter{smallq}[subq]

\newcommand{\SmallQ}{%
  \stepcounter{smallq}%
  \noindent\arabic{smallq}.\quad
}

\newcommand{\resetsmallq}{\setcounter{smallq}{0}}

\newcommand{\choices}[4]{\par
\noindent\hspace*{2em}A. #1\quad B. #2\quad C. #3\quad D. #4\par}

\newcommand{\choicelong}[4]{\par
\noindent\hspace*{2em}A. #1\par\hspace*{2em}B. #2\par\hspace*{2em}C. #3\par\hspace*{2em}D. #4\par}

\newcommand{\choicelongfive}[5]{\par
\noindent\hspace*{2em}A. #1\par\hspace*{2em}B. #2\par\hspace*{2em}C. #3\par\hspace*{2em}D. #4\par\hspace*{2em}E. #5\par}

% 分隔线
\newcommand{\examrule}{\noindent\rule{\linewidth}{0.6pt}}

% ---- 文档开始 ----
\begin{document}

% ============================================================
%  封面
% ============================================================
\begin{center}
  {\LARGE\bfseries 习近平新时代中国特色社会主义思想概论}\\[8pt]
  {\large\bfseries 全真模拟试卷（一）}\\[16pt]
  \begin{tabular}{lll}
    考试时间：120 分钟 & \quad & 满分：100 分 \\[4pt]
    \multicolumn{3}{l}{\textbf{注意事项：}请将答案写在答题纸指定区域，写在本卷上无效。}
  \end{tabular}
  \vspace{4pt}

  \examrule

  \vspace{6pt}
  \begin{tabular}{|c|c|c|c|c|}
    \hline
    \textbf{题号} & 一 & 二 & 三 & \textbf{总分} \\
    \hline
    \textbf{满分} & 30 & 20 & 50 & 100 \\
    \hline
    \textbf{得分} & & & & \\
    \hline
  \end{tabular}
\end{center}

\vspace{10pt}
\examrule
\vspace{6pt}

% ============================================================
%  第一大题：单项选择题（30题，每题1分，共30分）
% ============================================================
\BigQ{单项选择题（每题1分，共30分。在每小题给出的四个选项中，只有一项最符合题意。）}

\resetsmallq
\SmallQ 习近平新时代中国特色社会主义思想创立的时代背景，不包括以下哪一项？
\choices{世界百年未有之大变局加速演进}{中华民族伟大复兴进入关键时期}{中国特色社会主义进入新时代}{科学社会主义在21世纪的中国焕发新的蓬勃生机}
\vspace{4pt}
\SmallQ “两个确立”的决定性意义在于，它们是新时代取得一切成就的根本原因，是战胜一切艰难险阻、应对一切不确定性的最大确定性、最大底气、最大保证。其中“两个确立”是指：
\choicelong{确立习近平同志党中央的核心、全党的核心地位；确立习近平新时代中国特色社会主义思想的指导地位}{确立党中央的核心地位；确立马克思主义的指导地位}{确立习近平同志为党和国家最高领导人；确立中国特色社会主义道路}{确立党的全面领导；确立习近平新时代中国特色社会主义思想}
\vspace{4pt}
\SmallQ 中国特色社会主义进入新时代，是我国发展新的历史方位。以下关于“新时代”内涵的表述，错误的是：
\choicelong{是承前启后、继往开来、在新的历史条件下继续夺取中国特色社会主义伟大胜利的时代}{是决胜全面建成小康社会、进而全面建设社会主义现代化强国的时代}{是全国各族人民团结奋斗、不断创造美好生活、逐步实现全体人民同步富裕的时代}{是全体中华儿女勠力同心、奋力实现中华民族伟大复兴中国梦的时代}
\vspace{4pt}
\SmallQ 中国创造的人类文明新形态，以下表述不正确的是：
\choices{提供了全新的现代化路径，超越西方现代化模式}{为发展中国家提供了新的选择}{推进中国式现代化必须牢牢把握重大原则}{中国式现代化是对西方现代化模式的完全复制}
\vspace{4pt}
\SmallQ 推进中国式现代化必须牢牢把握的重大原则，下列哪项不属于五大原则？
\choices{坚持和加强党的全面领导}{坚持中国特色社会主义道路}{坚持全面依法治国}{坚持以人民为中心的发展思想}
\vspace{4pt}
\SmallQ “中国最大的国情就是中国共产党的领导”出自哪一章的核心考点？
\choices{第一章}{第二章}{第三章}{第四章}
\vspace{4pt}
\SmallQ 人民立场是中国共产党的根本政治立场，以下关于人民立场的表述，正确的是：
\choicelong{人民立场是党的唯一立场}{坚持人民立场，就要始终牢记党的初心和使命，始终保持党同人民群众的血肉联系}{人民立场是改革开放后才确立的}{人民立场与党的阶级立场是矛盾的}
\vspace{4pt}
\SmallQ 全面深化改革总目标是：
\choicelong{完善和发展中国特色社会主义制度，推进国家治理体系和治理能力现代化}{建设社会主义现代化强国}{实现中华民族伟大复兴}{全面建成小康社会}
\vspace{4pt}
\SmallQ 新发展理念的科学内涵中，创新、协调、绿色、开放、共享分别解决什么问题？
\choicelong{创新解决发展不平衡问题；协调解决发展动力问题；绿色解决人与自然和谐共生问题；开放解决发展内外联动问题；共享解决社会公平正义问题}{创新解决发展动力问题；协调解决发展不平衡问题；绿色解决人与自然和谐共生问题；开放解决发展内外联动问题；共享解决社会公平正义问题}{创新解决发展动力问题；协调解决社会公平正义问题；绿色解决发展不平衡问题；开放解决发展内外联动问题；共享解决人与自然和谐共生问题}{创新解决发展内外联动问题；协调解决发展动力问题；绿色解决社会公平正义问题；开放解决发展不平衡问题；共享解决人与自然和谐共生问题}
\vspace{4pt}
\SmallQ “科技强则国家强”强调了科技、人才、创新三者的关系，其中“第一动力”是指：
\choices{科技}{人才}{创新}{教育}
\vspace{4pt}
\SmallQ 人民民主是社会主义的生命，坚定不移走中国特色社会主义政治发展道路，其根本保证是：
\choices{党的领导}{人民当家作主}{依法治国}{民主集中制}
\vspace{4pt}
\SmallQ 文化繁荣兴盛是实现中华民族伟大复兴的必然要求，其中文化自信是更基础、更广泛、更深沉的力量。文化自信的深厚基础是：
\choices{革命文化}{社会主义先进文化}{中华优秀传统文化}{中国特色社会主义伟大实践}
\vspace{4pt}
\SmallQ 完善分配制度是推动高质量发展、促进社会公平的重要举措。分配制度是促进共同富裕的：
\choices{基础性制度}{根本性制度}{辅助性制度}{临时性措施}
\vspace{4pt}
\SmallQ 生态文明建设摆在全局工作的突出位置。以下哪项不是习近平生态文明思想的核心原则？
\choices{绿水青山就是金山银山}{保护生态环境就是保护生产力}{先污染后治理}{山水林田湖草沙一体化保护和系统治理}
\vspace{4pt}
\SmallQ “两个结合”中，马克思主义基本原理是“魂脉”，中华优秀传统文化是“根脉”。下列哪项正确描述了两者的关系？
\choices{两者是相互对立的}{两者是相互补充、相互融合的关系}{中华优秀传统文化是“魂脉”，马克思主义是“根脉”}{两者没有关系}
\vspace{4pt}
\SmallQ 习近平新时代中国特色社会主义思想的历史地位，以下表述正确的是：
\choicelong{是当代中国马克思主义、21世纪马克思主义，实现了马克思主义中国化时代化新的飞跃}{是当代中国马克思主义，但尚未形成新的飞跃}{是中国特色社会主义理论体系的重要组成部分，没有独立地位}{是中华文化和中国精神的时代精华，但不是马克思主义}
\vspace{4pt}
\SmallQ 我国社会主要矛盾已经转化为人民日益增长的美好生活需要和不平衡不充分的发展之间的矛盾。但“两个没有变”是指：
\choicelong{我国仍处于并将长期处于社会主义初级阶段的基本国情没有变；我国是世界最大发展中国家的国际地位没有变}{我国社会主义性质没有变；党的领导地位没有变}{我国基本经济制度没有变；分配制度没有变}{我国发展目标没有变；发展战略没有变}
\vspace{4pt}
\SmallQ 党的基本路线是国家的生命线、人民的幸福线，可概括为“一个中心、两个基本点”。其中“一个中心”是指：
\choices{以经济建设为中心}{以政治建设为中心}{以文化建设为中心}{以社会建设为中心}
\vspace{4pt}
\SmallQ 中国梦的本质是：
\choices{国家富强、民族振兴、人民幸福}{实现现代化}{建成小康社会}{中华民族伟大复兴}
\vspace{4pt}
\SmallQ “五个必由之路”中，团结奋斗是中国人民创造历史伟业的必由之路。团结奋斗最显著的精神标识是：
\choices{中国共产党}{中国人民}{中华民族}{中国特色社会主义}
\vspace{4pt}
\SmallQ 党的自身优势包括政治领导力、思想引领力、群众组织力、社会号召力。其中第一位的能力是：
\choices{思想引领力}{群众组织力}{社会号召力}{政治领导力}
\vspace{4pt}
\SmallQ 党中央集中统一领导是党的领导的最高原则。维护党中央权威和集中统一领导，是马克思主义执政党成熟的：
\choices{重大建党原则}{基本经验}{根本方法}{核心任务}
\vspace{4pt}
\SmallQ 民心是最大的政治，决定事业兴衰成败。群众路线是党始终坚守的根本工作方法，其中重要工作方法是：
\choices{典型引路}{调查研究}{以点带面}{群众运动}
\vspace{4pt}
\SmallQ 共同富裕是中国特色社会主义的本质要求，是中国式现代化的重要特征。推进共同富裕要坚持的原则是：
\choices{同步富裕、同等富裕}{有底线、渐进式推进}{平均主义}{先富后富齐头并进}
\vspace{4pt}
\SmallQ 社会主义基本经济制度的新概括包括公有制为主体、多种所有制经济共同发展，按劳分配为主体、多种分配方式并存，以及：
\choices{社会主义市场经济体制}{计划经济体制}{混合所有制经济}{股份合作制}
\vspace{4pt}
\SmallQ 新发展格局的核心是：
\choices{以国内大循环为主体、国内国际双循环相互促进}{国际循环为主}{内外循环均衡发展}{出口导向型经济}
\vspace{4pt}
\SmallQ 全过程人民民主体现“四个相统一”，以下哪项不属于这“四个相统一”？
\choices{过程民主和成果民主的统一}{程序民主和实质民主的统一}{直接民主和间接民主的统一}{选举民主和协商民主的统一}
\vspace{4pt}
\SmallQ 中国特色社会主义文化的三个主要组成部分是：中华优秀传统文化、革命文化和：
\choices{社会主义先进文化}{西方文化}{民俗文化}{大众文化}
\vspace{4pt}
\SmallQ 意识形态关乎旗帜、关乎道路、关乎国家安全。坚持马克思主义在意识形态领域指导地位的根本制度，是坚持和加强党对宣传文化事业全面领导的：
\choices{本质要求}{核心任务}{重要方针}{基本方法}
\vspace{4pt}
\SmallQ 伟大建党精神的内涵包括：坚持真理、坚守理想，践行初心、担当使命，不怕牺牲、英勇斗争，对党忠诚、不负人民。以下哪项是伟大建党精神的核心？
\choices{理想信念}{为民服务}{艰苦奋斗}{忠诚担当}

\vspace{10pt}
\examrule

% ============================================================
%  第二大题：多项选择题（10题，每题2分，共20分）
% ============================================================
\BigQ{多项选择题（每题2分，共20分。每小题有两个或两个以上符合题意，多选、少选、错选均不得分。）}

\resetsmallq
\SmallQ 习近平新时代中国特色社会主义思想是完整的科学体系，其组成部分包括：
\choices{“十个明确”}{“十四个坚持”}{“十三个方面成就”}{“六个必须坚持”}
\vspace{4pt}
\SmallQ 推进中国式现代化需要正确处理的重大关系包括：
\choices{顶层设计与实践探索的关系}{战略与策略的关系}{守正与创新的关系}{效率与公平的关系}
\vspace{4pt}
\SmallQ 关于“四个自信”，以下说法正确的是：
\choicelong{道路自信是对发展方向和未来命运的自信}{理论自信是对马克思主义理论特别是中国特色社会主义理论体系的科学性、真理性的自信}{制度自信是对中国特色社会主义制度具有制度优势的自信}{文化自信是更基础、更广泛、更深厚的自信}
\vspace{4pt}
\SmallQ 全过程人民民主的特征包括：
\choices{全链条的民主}{全方位的民主}{全覆盖的民主}{最广泛、最真实、最管用的民主}
\vspace{4pt}
\SmallQ 全面深化改革开放坚持正确方法论，具体要求包括：
\choices{增强系统性、整体性、协同性}{加强顶层设计和摸着石头过河相结合}{统筹改革发展稳定}{胆子要大，步子要稳}
\vspace{4pt}
\SmallQ 新发展理念的科学内涵包括：
\choicelong{创新是引领发展的第一动力}{协调是持续健康发展的内在要求}{绿色是永续发展的必要条件和人民对美好生活追求的重要体现}{开放是国家繁荣发展的必由之路，共享是中国特色社会主义的本质要求}
\vspace{4pt}
\SmallQ 关于“两个确立”和“两个维护”，以下表述正确的有：
\choicelong{“两个确立”是新时代取得一切成就的根本原因}{“两个维护”中关键是维护党中央权威和集中统一领导}{“两个确立”与“两个维护”内涵完全一样，可以互换使用}{两者紧密联系但内涵不同，应结合起来复习}
\vspace{4pt}
\SmallQ 以下关于中国特色社会主义政治发展道路的表述，正确的有：
\choices{必须坚持党的领导、人民当家作主、依法治国有机统一}{党的领导是人民当家作主和依法治国的根本保证}{人民当家作主是社会主义民主政治的本质特征}{依法治国是党领导人民治理国家的基本方式}
\vspace{4pt}
\SmallQ 文化建设中，以下属于中国特色社会主义文化组成部分的有：
\choices{中华优秀传统文化}{革命文化}{社会主义先进文化}{西方现代文化}
\vspace{4pt}
\SmallQ 关于共同富裕，以下理解正确的是：
\choices{共同富裕是中国特色社会主义的本质要求}{共同富裕是中国式现代化的重要特征}{共同富裕要求所有人同步富裕、同等富裕}{推进共同富裕要坚持有底线、渐进式推进}

\vspace{10pt}
\examrule

% ============================================================
%  第三大题：简答题（5题，每题10分，共50分）
% ============================================================
\BigQ{简答题（每题10分，共50分。要求简明扼要地回答问题，并展开简要论述，只写结论不得满分。）}

\resetsmallq
\SmallQ（10分）试述习近平新时代中国特色社会主义思想创立的时代背景（五个方面）。
\vspace{6pt}
\SmallQ（10分）如何理解“中国最大的国情是中国共产党的领导”？请结合党的领导制度优势简要论述。
\vspace{6pt}
\SmallQ（10分）推进中国式现代化必须牢牢把握的重大原则有哪些？请分别简要说明。
\vspace{6pt}
\SmallQ（10分）简述新发展理念的科学内涵及其各自要解决的问题。
\vspace{6pt}
\SmallQ（10分）为什么说“文化繁荣兴盛是实现中华民族伟大复兴的必然要求”？请从多个角度加以论述。

\vspace{16pt}
\examrule
\vspace{4pt}
\begin{center}
  {\large\bfseries ——试题到此结束，请认真检查——}
\end{center}
\vspace{4pt}
\examrule

% ============================================================
%  参考答案
% ============================================================
\newpage
\setcounter{bigq}{0}

\begin{center}
  {\LARGE\bfseries 参考答案与评分标准}
\end{center}
\vspace{8pt}
\examrule

% -------- 第一大题参考答案 --------
\BigQ{单项选择题（每题1分，共30分）}

\begin{answerbox}
\begin{enumerate}[label=\arabic*., itemsep=0pt]
  \item C（中国特色社会主义进入新时代是历史方位，不是时代背景）
  \item A
  \item C（应为“逐步实现全体人民共同富裕”，不是“同步富裕”）
  \item D（中国式现代化不是复制西方模式，而是创造新形态）
  \item C（依法治国不属于五大原则，五大原则为：党的领导、中国道路、人民中心、改革开放、斗争精神）
  \item C
  \item B
  \item A
  \item B
  \item C（创新是第一动力）
  \item A
  \item C
  \item A
  \item C
  \item B
  \item A
  \item A
  \item A
  \item A
  \item B（团结奋斗是中国人民最显著的精神标识）
  \item D
  \item A
  \item B
  \item B
  \item A
  \item A
  \item D（应为“全过程人民民主”体现过程与成果、程序与实质、直接与间接、人民民主与国家意志四统一，本题D项“选举民主和协商民主的统一”不属于该四个统一）
  \item A
  \item A
  \item A（伟大建党精神的核心是理想信念，即“坚持真理、坚守理想”）
\end{enumerate}
\end{answerbox}

% -------- 第二大题参考答案 --------
\BigQ{多项选择题（每题2分，共20分）}

\begin{answerbox}
\begin{enumerate}[label=\arabic*., itemsep=0pt]
  \item ABCD（“十个明确”“十四个坚持”“十三个方面成就”“六个必须坚持”共同构成科学体系）
  \item ABCD（六对重大关系包含以上四项，另两项为活力与秩序、自立自强与对外开放）
  \item ABCD
  \item ABCD
  \item ABCD
  \item ABCD
  \item ABD（C错误，两者内涵不同，不能混用）
  \item ABCD
  \item ABC
  \item ABD（C错误，共同富裕不是同步同等富裕，而是渐进式推进）
\end{enumerate}
\end{answerbox}

% -------- 第三大题参考答案 --------
\BigQ{简答题（每题10分，共50分）}

\begin{answerbox}

\textbf{（1）习近平新时代中国特色社会主义思想创立的时代背景（五个方面）}

答：①世界百年未有之大变局加速演进；②中华民族伟大复兴进入关键时期；③中国式现代化全面推进拓展；④科学社会主义在21世纪的中国焕发新的蓬勃生机；⑤中国共产党在革命性锻造中更加坚强有力。这五个方面共同构成了新思想产生的时代条件，体现了理论与实践、国内与国际、历史与现实的统一。

\textbf{（2）如何理解“中国最大的国情是中国共产党的领导”？}

答：这一论断深刻揭示了党的领导与中国特色社会主义的内在关系。第一，中国共产党领导是中国特色社会主义最本质的特征，是中国特色社会主义制度的最大优势；第二，党的领导是做好党和国家各项工作的根本保证，是战胜一切困难和风险的“定海神针”；第三，党的自身优势（政治领导力、思想引领力、群众组织力、社会号召力）是中国特色社会主义制度优势的主要来源。因此，中国最大的国情就是中国共产党的领导，这是从历史和现实、理论和实践结合上得出的科学结论。

\textbf{（3）推进中国式现代化必须牢牢把握的重大原则}

答：五大原则为：①坚持和加强党的全面领导——确保现代化建设的正确方向；②坚持中国特色社会主义道路——不走封闭僵化的老路和改旗易帜的邪路；③坚持以人民为中心的发展思想——现代化成果由人民共享；④坚持深化改革开放——为现代化注入强大动力；⑤坚持发扬斗争精神——应对风险挑战，开创事业新局面。这五条原则相互联系、相互支撑，构成推进中国式现代化的根本遵循。

\textbf{（4）新发展理念的科学内涵及其各自要解决的问题}

答：①创新是引领发展的第一动力，解决发展动力问题；②协调是持续健康发展的内在要求，解决发展不平衡问题；③绿色是永续发展的必要条件，解决人与自然和谐共生问题；④开放是国家繁荣发展的必由之路，解决发展内外联动问题；⑤共享是中国特色社会主义的本质要求，解决社会公平正义问题。五大理念相互贯通，完整、准确、全面贯彻新发展理念是关系我国发展全局的一场深刻变革。

\textbf{（5）为什么说“文化繁荣兴盛是实现中华民族伟大复兴的必然要求”？}

答：可以从四个维度论述：第一，文化繁荣兴盛是实现中华民族伟大复兴的精神支撑，文化兴则国运兴；第二，是建设社会主义现代化强国的应有之义，现代化强国必然包括文化强国；第三，是满足人民日益增长的美好生活需要的内在要求，人民对精神文化生活的需求日益增长；第四，是在世界文化激荡中站稳脚跟的前提基础，文化自信是更基础、更广泛、更深厚的自信。因此，文化繁荣兴盛直接关系到民族复兴全局，必须坚持中国特色社会主义文化发展道路，激发全民族文化创新创造活力。

\end{answerbox}

\vspace{12pt}
\examrule
\begin{center}
  \textit{参考答案仅供参考，简答题需结合教材理论展开论述，论述